\documentclass[11pt,a4paper]{article}
\usepackage[italian]{babel}
\usepackage[margin=2.5cm]{geometry}
\usepackage{parskip}
\usepackage{amsmath, amssymb, amsthm}
\usepackage{booktabs}
\usepackage{array}
\usepackage{xcolor}
\usepackage{needspace}
\usepackage{enumitem}
\usepackage{listings}
\usepackage{tikz}
\usepackage[most]{tcolorbox}
\usepackage[hidelinks]{hyperref}
\usetikzlibrary{decorations.pathreplacing, shapes.geometric, arrows.meta, positioning}

% Niente vedove/orfane: i paragrafi non lasciano righe isolate a fine/inizio pagina
\widowpenalty=10000
\clubpenalty=10000

% Elenchi più compatti
\setlist{itemsep=2pt, parsep=0pt, topsep=4pt, partopsep=0pt}

% --- Riquadri con bordo nero, senza numero (si spezzano tra due pagine) ---
% Uso: \begin{definizione} ... \end{definizione}
%      \begin{teorema}[Nome] ... \end{teorema}  (il nome è facoltativo)
\tcbset{nota/.style={enhanced jigsaw, breakable, colback=white,
  colframe=black, boxrule=0.5pt, sharp corners}}
\NewTColorBox{definizione}{}{nota,
  before upper={\textbf{Definizione.}\ }}
\NewTColorBox{teorema}{o}{nota,
  before upper={\textbf{Teorema\IfValueT{#1}{ (#1)}.}\ }}
\NewTColorBox{proposizione}{o}{nota,
  before upper={\textbf{Proposizione\IfValueT{#1}{ (#1)}.}\ }}
\NewTColorBox{proprieta}{o}{nota,
  before upper={\textbf{Proprietà\IfValueT{#1}{ (#1)}.}\ }}
\NewTColorBox{assioma}{o}{nota, boxrule=1pt,
  before upper={\textbf{Assioma\IfValueT{#1}{ (#1)}.}\ }}

% --- Ambienti semplici (senza riquadro) ---
% Il titolo sta in cima al blocco, anche se il contenuto inizia con un riquadro o
% con codice; e non resta mai da solo in fondo a una pagina.
% Uso: \begin{esempio}[Titolo facoltativo] ... \end{esempio}
\newcommand{\newrunin}[2]{%
  \NewDocumentEnvironment{#1}{o}{%
    \par\Needspace{4\baselineskip}\addvspace{\medskipamount}\noindent
    \textbf{#2}\IfValueT{##1}{ (##1)}\textbf{.}\ \ignorespaces}%
  {\par\addvspace{\medskipamount}}}
\newrunin{esempio}{Esempio}
\newrunin{esercizio}{Esercizio}
\newrunin{osservazione}{Osservazione}

% V e F nelle tavole di verità
\newcommand{\V}{\textsf{V}}
\newcommand{\F}{\textsf{F}}

% --- Codice C ---
% Uso: \begin{codice} ... \end{codice}      codice C numerato
%      \begin{comando} ... \end{comando}    comandi da terminale
%      \begin{risultato} ... \end{risultato} output di un programma
%      \lstinline|printf("ciao");|           codice dentro al testo
\lstdefinestyle{c}{
  language=C,
  basicstyle=\ttfamily\small,
  keywordstyle=\color{blue}\bfseries,
  commentstyle=\color{gray}\itshape,
  stringstyle=\color{red!70!black},
  numbers=left,
  numberstyle=\tiny\color{gray},
  numbersep=8pt,
  columns=flexible,
  keepspaces=true,
  breaklines=true,
  showstringspaces=false,
  tabsize=4,
  morekeywords={include, define},
  % lettere accentate nei commenti
  literate={à}{{\`a}}1 {è}{{\`e}}1 {é}{{\'e}}1 {ì}{{\`i}}1
           {ò}{{\`o}}1 {ù}{{\`u}}1 {È}{{\`E}}1 {À}{{\`A}}1
}
\lstset{style=c}
\newtcblisting{codice}{listing only, nota, left=6mm,
  listing options={style=c}}
\newtcblisting{comando}{listing only, nota,
  listing options={basicstyle=\ttfamily\small, numbers=none, language={},
    columns=flexible, keepspaces=true, breaklines=true}}
\newtcblisting{risultato}{listing only, enhanced jigsaw, breakable,
  colback=black!5, colframe=black, boxrule=0.5pt, sharp corners,
  listing options={basicstyle=\ttfamily\small, numbers=none, language={},
    columns=flexible, keepspaces=true, breaklines=true}}

% --- Diagrammi di flusso (simboli standard) ---
\tikzset{
  terminale/.style = {ellipse, draw, minimum width=2.5cm, minimum height=0.9cm},
  io/.style        = {trapezium, trapezium left angle=70, trapezium right angle=110,
                      draw, minimum width=2.5cm, minimum height=0.9cm},
  processo/.style  = {rectangle, draw, minimum width=2.5cm, minimum height=0.9cm},
  decisione/.style = {diamond, draw, aspect=2, minimum width=2.5cm},
  freccia/.style   = {thick, -{Stealth}}
}

\title{Il teorema di sostituzione}
\author{Marco Cerbella}
\date{Lezione 14 -- 10 novembre 2026}

\begin{document}
\maketitle

\section{Teorema di sostituzione}

\begin{teorema}[Teorema di sostituzione]
Supponiamo che esista (finito o infinito)
\[
  \lim_{x\to c} f(x) = \ell .
\]
Sia poi $g$ una funzione definita in un intorno di $\ell$ (escluso al più il punto $\ell$) e tale che:
\begin{itemize}
  \item[(i)] se $\ell \in \mathbb{R}$, $g$ è continua in $\ell$;
  \item[(ii)] se $\ell = +\infty$ oppure $\ell = -\infty$, esiste finito o infinito $\displaystyle\lim_{y\to\ell} g(y)$.
\end{itemize}
Allora esiste il limite per $x$ tendente a $c$ della funzione composta $g\circ f$ e si ha
\[
  \lim_{x\to c} g(f(x)) = \lim_{y\to\ell} g(y).
\]
\end{teorema}

Nel caso (i) si ottiene
\[
  \lim_{x\to c} g(f(x)) = g\!\left(\lim_{x\to c} f(x)\right).
\]

\begin{osservazione}[Variante (i$'$)]
Se $\ell \in \mathbb{R}$ e esiste un intorno $I(c)$ di $c$ in cui $f(x) \neq \ell$ per ogni $x \neq c$, e esiste finito o infinito $\displaystyle\lim_{y\to\ell} g(y)$, allora vale ancora $\displaystyle\lim_{x\to c} g(f(x)) = \lim_{y\to\ell} g(y)$ (senza richiedere che $g$ sia continua in $\ell$, né definita in $\ell$).
\end{osservazione}

\begin{esempio}[$\lim_{x\to 0}\left(\frac{\sin x}{x}\right)^{2}$]
Si pone $f(x) = \dfrac{\sin x}{x}$, $\ell = 1$ e $g(y) = y^2$. Si ha $\displaystyle\lim_{x\to 0}\frac{\sin x}{x} = 1$ e $g$ è continua in $\ell = 1$: siamo nelle ipotesi (i) del teorema di sostituzione, quindi
\[
  \lim_{x\to 0}\left(\frac{\sin x}{x}\right)^{2} = \lim_{y\to 1} g(y) = \lim_{y\to 1} y^2 = 1,
\]
cioè $\displaystyle\lim_{x\to 0}\left(\frac{\sin x}{x}\right)^{2} = \left(\lim_{x\to 0}\frac{\sin x}{x}\right)^{2} = 1^2 = 1$.
\end{esempio}

\begin{esempio}[$\lim_{x\to 0}\frac{\sin(x^2)}{x^2}$]
Si pone $y = x^2$: $f(x) = x^2$, $\displaystyle\lim_{x\to 0} x^2 = 0 = \ell$. Si considera
\[
  g(y) = \begin{cases} \dfrac{\sin y}{y} & y \neq 0 \\[2mm] 1 & y = 0 \end{cases}
\]
(estensione con continuità della funzione $\frac{\sin y}{y}$): infatti $\displaystyle\lim_{y\to 0} g(y) = \lim_{y\to 0}\frac{\sin y}{y} = 1 = g(0)$, quindi $g$ è continua in $0$. Siamo nel caso (i) del teorema di sostituzione:
\[
  \lim_{x\to 0}\frac{\sin(x^2)}{x^2} = \lim_{y\to 0} g(y) = \lim_{y\to 0}\frac{\sin y}{y} = 1.
\]
\end{esempio}

\begin{esempio}[$\lim_{x\to 1}\arctan\frac{1}{x-1}$]
Sia $h(x) = \arctan\dfrac{1}{x - 1}$. Il limite $\displaystyle\lim_{x\to 1}\frac{1}{x-1}$ non esiste, perché
\[
  \lim_{x\to 1^+}\frac{1}{x-1} = +\infty, \qquad \lim_{x\to 1^-}\frac{1}{x-1} = -\infty .
\]
Si calcolano allora i limiti destro e sinistro con il teorema di sostituzione, con $f(x) = \dfrac{1}{x-1}$ e $g(y) = \arctan y$ (funzione continua, con limiti agli estremi $\pm\dfrac{\pi}{2}$). Posto $y = \dfrac{1}{x-1}$:
\[
  \lim_{x\to 1^+}\arctan\frac{1}{x - 1} = \lim_{y\to+\infty}\arctan y = \frac{\pi}{2}, \qquad
  \lim_{x\to 1^-}\arctan\frac{1}{x - 1} = \lim_{y\to-\infty}\arctan y = -\frac{\pi}{2}.
\]
I limiti destro e sinistro sono finiti ma diversi: il limite bilatero non esiste e $x = 1$ è un punto di \emph{discontinuità a salto}. Il salto vale
\[
  \lim_{x\to 1^+}h(x) - \lim_{x\to 1^-}h(x) = \frac{\pi}{2} - \left(-\frac{\pi}{2}\right) = \pi .
\]
\end{esempio}

\begin{esempio}[$\lim_{x\to+\infty}\log\left(\sin\frac{1}{x}\right)$]
Si pone $y = \dfrac{1}{x}$. Poiché $\displaystyle\lim_{x\to+\infty}\frac{1}{x} = 0$ e $\sin y$ è una funzione continua in $y = 0$,
\[
  \lim_{x\to+\infty}\sin\frac{1}{x} = \lim_{y\to 0}\sin y = 0 .
\]
Sia $z = \sin\dfrac{1}{x}$: $\displaystyle\lim_{z\to 0^+}\log z = -\infty$. Il limite di $z$ è $0^+$, cioè $z \to 0$ con $z > 0$, perché $\sin\dfrac{1}{x} > 0$ per $x > \dfrac{1}{\pi}$, cioè quando $\dfrac{1}{x} < \pi$ e $\dfrac{1}{x} > 0$. Quindi
\[
  \lim_{x\to+\infty}\log\left(\sin\frac{1}{x}\right) = \lim_{z\to 0^+}\log z = -\infty .
\]
\end{esempio}

\end{document}
